\scanpage{95}
\mspage{101}

\uline{Définition 1} $f:X\to Y$ dans $\mathrm{Diag}$ est appelé $\mathbb D$-\uline{équivalence cohomologique} (sous-entendu : \uline{complète}) si pour tout $F\deleted{\unsure{\in\mathrm{Ob}|\in}}\mathbb D(Y)$, $F$ \uline{loc. constant}, on a

\[
H^{\cdot}_{\mathbb D}(Y,F)\xrightarrow{\sim} H^{\cdot}_{\mathbb D}(X,\unsure{f|f^*}(F))
\]

Notion duale $\mathbb D$-\uline{équivalence homologique}

\[
H^{\mathbb D}_{\cdot}(X,\unsure{f^*|f}(F))\xrightarrow{\sim}H^{\mathbb D}_{\cdot}(Y,F).
\]

Si $Y$ est \deleted{\unsure{total|tota}}\unsure{plètement|complètement} $\mathbb D$-asphérique, alors les coeff loc constants \inserted{sur $Y$} sont constants, et $f$ est une $\mathbb D$-équiv. coh. \inserted{(ou hom.)} \unsure{ssi|si} c'est une $\mathbb D$-équivalence, i.e. si $X$ est $\mathbb D$-asphérique.

% margin note written sideways
\margin{L'axiome \uline{Der} 3 b'$_1$ implique que toute $\mathrm{Hot}$-équivalence est une $\mathbb D$-équiv. cohom.}

\uline{Prop. 1} \inserted{(sorites)} : Soient $X\xrightarrow{f}Y\xrightarrow{g}Z$. Si $f$ et $g$ sont des $\mathbb D$-équivalences coh, alors $gf$ aussi ; si $gf$ et $f$ sont des $\mathbb D$-équiv. coh, alors $g$ aussi. Si $gf$ est une $\mathbb D$-équiv$^{\mathrm{ce}}$ coh, et si $g$ est une $\mathbb D$-équiv. \uline{2-complète} (donc tout $F$ loc constant sur $Y$ provient d'un $G$ loc constant sur $Z$), alors $f$ est une $\mathbb D$-équiv. coh.

\uline{Définition 2} : 1°) $f:X\to Y$ dans $\mathrm{Diag}$ est \deleted{appelé} une $\mathbb D$-\uline{fibration cohomologique} \inserted{(\deleted{\unsure{stricte|absolue}}\uuline{complète})}, je dis bien \uuline{complète} si

a) $\forall\,y\in\mathrm{Ob}\,Y$, $X_y\unsure{\relbar|\to}X_{/y}$ est une $\mathbb D$-équiv. \deleted{1-complète} cohomologique (complète)

b) $\forall\,y'\to y$ dans $Y$, $X_{/y'}\unsure{\relbar|\to}X_y$ est une $\mathbb D$-équiv. \deleted{1 complète} cohomologique (complète)

% archive sheet number: 94
